\section{Exact order-constrained depolarization in arbitrary dimension}\label{sec:order}
Let \(\E_A(X)=I_d/d\otimes\Tr_AX\) and
\(\N_\lambda=\lambda\id+(1-\lambda)\E_A\), where
\(-1/(d^2-1)\le\lambda\le1\), the known physical range~\cite[Section~8.1]{HRF}. For \(\kappa>1\) restrict the finite pairs to
\begin{equation}\label{eq:order-class}
 \rho_R=\sigma_R,\quad \rho,\sigma>0,\quad \rho\ne\sigma,\quad
 \rho\preceq\kappa\E_A(\rho),\quad \sigma\preceq\kappa\E_A(\sigma).
\end{equation}
For the metric impose the base order bound and \(\Tr_AD=0\); one may also require a two-sided small neighborhood satisfying \eqref{eq:order-class}. The suprema proved below are unchanged. The budget uses the \emph{actual} reference marginal, not an optimized conditional min-entropy reference.

\begin{proposition}[Commuting Bell lower and its range]\label{prop:bell}
For every \(d\ge2\), \(1<\kappa\le d^2/2\) and every nonaffine operator-convex Petz generator normalized by \(f(1)=0\), the complete pair coefficient under \eqref{eq:order-class} is
\(\lambda^2\kappa/[1+\lambda(\kappa-1)]\) throughout the physical depolarizing range. The SLD and minimal-chi-square coefficients have the same value. A reference of dimension \(d\) suffices for the lower supremum.
\end{proposition}
\begin{proof}
The order upper is proved below. Write \(n=d^2\), \(B=I_n/n\), and take two orthogonal maximally entangled Bell projectors \(P_0,P_1\). For \(1<K<\kappa\le n/2\), set
\[
 a=K/n,\quad b=(1-2a)/(n-2),\quad
 \sigma_K=bI_n+(a-b)(P_0+P_1),\quad D=P_0-P_1.
\]
Here \(a>b>0\), \(\E_A(\sigma_K)=B\), \(\Tr_AD=0\), and the actual order constant is \(K\). The input and output commute; their eigenvalues on the two changed rays are \(a\) and \(a_\lambda=\lambda a+(1-\lambda)/n>0\). Hence the SLD ratio is exactly \(\lambda^2a/a_\lambda\).

For \(0<|t|<\min(a,b,(\kappa-K)/n)\), \(\rho_t=\sigma_K+tD\) and \(\sigma_K\) are faithful with identical reference and both order bounds. Their Petz divergences are
\[
 a\left[f(1+t/a)+f(1-t/a)\right],\qquad
 a_\lambda\left[f(1+\lambda t/a_\lambda)+f(1-\lambda t/a_\lambda)\right].
\]
For a normalized nonaffine operator-convex generator, \(f''(1)>0\), as follows from the nonzero positive representation in \eqref{eq:representation}. The Taylor ratio as \(t\to0\) is therefore \(\lambda^2a/a_\lambda=\lambda^2K/[1+\lambda(K-1)]\). Then let \(K\uparrow\kappa\). At the endpoint \(\kappa=n/2\), each earlier base remains faithful. The cases \(\lambda=0,1\) are immediate. The quadratic chi-square equality follows directly from the metric ratio.
\end{proof}

This is a reparametrization of the Bell family already present in our earlier component work, not an independent invention of Bell-assisted contraction. In particular, when \(d\ge3\), the interval \(\kappa\le d+1\) is contained in \(\kappa\le d^2/2\); the rank-\(d\) ladder previously considered on that interval is an alternative proof. A conditional-expectation index alone does not imply sharpness: the pinching \(\E=(\id+\Ad_Z)/2\) has index two, but its equal-\(\E\) contraction is \(\lambda^2\), rather than the larger raw order comparison, for \(0<\lambda<1\).

\begin{theorem}[Rank-two controlled order profile]\label{thm:order}
For \(d\ge2\) and \(1<\kappa\le L_d=d(d+1)/2\), the SLD, minimal-chi-square and finite triangular-Petz contraction coefficients under \eqref{eq:order-class} equal
\begin{equation}\label{eq:profile}
 C_\kappa(\lambda)=\frac{\lambda^2\kappa}{1+\lambda(\kappa-1)}.
\end{equation}
A reference of dimension \(d+1\) suffices for each lower supremum. The same expression is an upper for every nonaffine operator-convex Petz generator normalized by \(f(1)=0\), and is sharp as a bound uniform over that family. No exact higher-dimensional coefficient beyond the stated order range is asserted.
\end{theorem}

\begin{proof}
The standard order comparison gives, for each admissible state,
\[
 \N_\lambda(\sigma)\succeq c\sigma,\qquad c=\lambda+(1-\lambda)/\kappa>0.
\]
The positivity follows from \(1+\lambda(\kappa-1)\ge1-(L_d-1)/(d^2-1)>0\).
On equal-reference differences \(\N_\lambda(D)=\lambda D\). Apply inverse antitonicity to \((L_\sigma+R_\sigma)/2\), or to \(L_\rho+tR_\sigma\) using the order for both states. The quadratic endpoint and positive integral in \eqref{eq:resolvent} give the upper \(\lambda^2/c\). This comparison machinery is prior; the matching lower is the construction below.

Put \(m=d-1\), \(q=1+me^2\), and give \(R\) the labels \(0,\ldots,d\). With A indices modulo \(d\), set
\[
 a_i=|i,i\rangle+e\sum_{k=1}^{d-1}|i+k,k+1\rangle\quad(i=0,1),\qquad
 \sigma_e=\frac{a_0a_0^*+a_1a_1^*}{2q}.
\]
The two vectors have disjoint supports and squared norm \(q\), so \(P=2\sigma_e\) is a projection. The reference marginal has mass \(1/(2q)\) on labels zero and one, and mass \(e^2/q\) on each remaining label. Hence \(B_e=\E_A(\sigma_e)>0\).
The whitened normalized vectors are orthogonal and have squared norm
\(d+d(d-1)/2=L_d\); consequently \(B_e^{-1/2}\sigma_e B_e^{-1/2}\) has two eigenvalues \(L_d\) and all others zero.

Let \(b_0=\sum_{k=1}^{d-1}|k,k+1\rangle\), \(b_1=-\sum_{k=1}^{d-1}|1+k,k+1\rangle\) and
\(D=\sum_{i=0}^1(b_i a_i^*+a_i b_i^*)\). All off-diagonal partial-trace terms vanish because their A indices differ. At each center reference label the opposite diagonal derivatives cancel. Thus \(\Tr_AD=0\), \((I-P)D(I-P)=0\), and the finite boundary SLD form is \(16mq\).

For \(0\le\theta<1\), define \(s_{e,\theta}=(1-\theta)B_e+\theta\sigma_e\). Its whitened eigenvalues are \(1+(L_d-1)\theta\) twice and \(1-\theta\) otherwise. It is faithful and has exactly that order constant. Moreover
\(\N_\lambda(s_{e,\theta})=s_{e,\lambda\theta}\).

Each vector's \(d\)-dimensional support is an invariant block. On the main coordinate and the uniform center superposition its base and tangent reduce to
\[
 W_p=\frac1{2dq}\begin{pmatrix}
 1+(d-1)p&dpe\sqrt m\\dpe\sqrt m&[2+(dm-2)p]e^2
 \end{pmatrix},\quad
 T=\begin{pmatrix}0&\sqrt m\\\sqrt m&2em\end{pmatrix},
\]
with opposite signs for the two tangent blocks. The other \(m-1\) center modes have base eigenvalue \(e^2(1-p)/(dq)\) and zero tangent. All remaining coordinates are positive diagonal replacement contributions with zero tangent. The aligned determinant is
\[
 \det W_p=\frac{e^2(1-p)[1+(L_d-1)p]}{2d^2q^2}>0
\]
for the values \(p=\theta,\lambda\theta\) used here.

Let \(G_e(p)=\Tr(TY)\), where \(W_pY+YW_p=2T\). Scaling the second diagonal unknown by \(e\), or eliminating it directly, gives
\begin{equation}\label{eq:general-limit}
 \lim_{e\downarrow0}\frac{G_e(p)}{8m}=\frac{L_d}{1+(L_d-1)p}.
\end{equation}
For completeness, put \(A=1+(d-1)p\), \(B=dp\sqrt m\), \(C=2+(dm-2)p\). The limiting system for the off-diagonal unknown and the scaled diagonal has matrix \(\left(\begin{smallmatrix}A&B\\B&C\end{smallmatrix}\right)\). Its determinant is \(2(1-p)[1+(L_d-1)p]\); the numerator from the direction \((\sqrt m,m)\) is \(m[C-2dmp+Am]=m(m+2)(1-p)\). This proves \eqref{eq:general-limit} at every fixed \(-1/(L_d-1)<p<1\), without assuming uniform convergence at \(p=1\).

Both blocks have the same retention ratio. Therefore, at fixed \(\theta<1\), the ratio tends to
\[
 \lambda^2\frac{1+(L_d-1)\theta}{1+(L_d-1)\lambda\theta}.
\]
Choose \(\theta<(\kappa-1)/(L_d-1)\). Both \(s_{e,\theta}\) and \(\kappa B_e-s_{e,\theta}\) are strictly positive. For each finite \(e,\theta\), the pair \(\rho_t=s_{e,\theta}+tD\), \(\sigma=s_{e,\theta}\) stays faithful with both order bounds for sufficiently small nonzero real \(t\). Its reference marginal stays fixed. For triangular Petz take \(t\to0\) first; for minimal chi-square the ratio is exactly quadratic. Then take \(e\downarrow0\) and finally \(\theta\uparrow(\kappa-1)/(L_d-1)\). This proves \eqref{eq:profile}. The cases \(\lambda=0,1\) are the zero-on-differences and identity endpoints.
\end{proof}

\begin{remark}[The empty endpoint and the domain of the order bound]
At \(\kappa=1\), positivity and equal traces force \(\rho=\E_A(\rho)\) and \(\sigma=\E_A(\sigma)\); equal reference then forces equality. There are no distinct finite pairs, so \eqref{eq:profile} is not assigned to that empty optimization. The base-only SLD problem with unrestricted tangents at a product base is a different problem. Likewise the expression \(\lambda^2/c\) must never be used when \(c\le0\); for example \(\lambda=-1/3,\kappa=5\) gives a negative raw expression in dimension two, not a contraction bound.
\end{remark}

\section{Exact qubit resource restrictions}\label{sec:resources}
For \(d=2\), Theorems~\ref{thm:pauli} and~\ref{thm:order} combine to give, for every \(\kappa>1\),
\[
 C_\kappa(\lambda)=\frac{\lambda^2K}{1+\lambda(K-1)},\qquad K=\min(\kappa,3),\qquad -1/3\le\lambda\le1.
\]
For \(\kappa\ge3\), the upper is the unrestricted isotropic case of Theorem~\ref{thm:pauli}, \(3\lambda^2/(1+2\lambda)\); the same lower family has order arbitrarily close to three. At the channel endpoints use DPI or the identity directly.

For a qubit and any finite reference,
\[
 2\E_A(\sigma)-\sigma=(\sigma_y\otimes I)\sigma^{T_A}(\sigma_y\otimes I).
\]
Thus PPT bases obey the order-two bound, as do separable bases. The exact complete SLD/minimal-chi-square/triangular coefficient restricted to PPT or separable models is
\(2\lambda^2/(1+\lambda)\). For finite Petz both states must belong to the class.
To prove the matching lower, use the qubit ladder above with \(s_{e,\theta}\) and \(\theta<1/2\). Its partial transpose has two blocks with determinant
\(e^2(1-\theta-2\theta^2)/[8(1+e^2)^2]>0\) and two positive singletons. Small bilateral perturbations stay PPT. In dimension \(2\times3\), the prior Horodecki theorem~\cite{HHH} converts these pairs to separable ones. This uses PPT=SEP only for the lower in \(2\times3\), not for arbitrary references or larger local dimensions.

For \emph{each} normalized nonaffine operator-convex Petz generator, the SEP/PPT pair coefficient is also exactly \(2\lambda^2/(1+\lambda)\). Indeed, the order-two comparison gives the upper, while the qubit Bell family of Proposition~\ref{prop:bell}, with \(K<2\), has all Bell probabilities \(p_j\) strictly below \(1/2\). The partial-transpose eigenvalues are \(1/2-p_j\). Both the base and sufficiently small finite perturbations are strictly PPT, hence separable in \(2\times2\), and approach the stated lower. The same gives the budget profile inside SEP/PPT for \(1<\kappa\le2\). This restricted-pair statement includes Umegaki but does not determine the unrestricted coefficient or actual MI, whose second state has an additional product-marginal requirement.

A classical-reference \emph{family} has both states and the tangent block diagonal in one fixed reference basis, with fixed marginal weights. The forms split into weighted ordinary-qubit forms, whose coefficient is \(\lambda^2\) by~\cite{HR}; \(R=1\) gives equality. Classicality of the base alone is insufficient. For example at \(\lambda=-1/3\), the classical base \(\sigma=\operatorname{diag}(9,1,1,9)/20\) with \(D_{03}=D_{30}=1\) gives retention \(3/11>1/9\); the direction creates reference coherences.

At \(\lambda=-1/3\) the exact levels are therefore
\[
 \text{classical-reference family}:\frac19,\qquad
 \text{SEP/PPT family}:\frac13,\qquad
 \text{unrestricted quantum family}:1.
\]
SLD retention (or minimal-chi-square retention with base \(\sigma\)) above \(2\lambda^2/(1+\lambda)\) is a sufficient certificate that that input base is NPT. For a general finite Petz pair, the proved implication is only that exceeding this threshold excludes both input states being PPT; it does not identify which state is NPT. This mathematical implication is a contrapositive of the prior reduction criterion plus the quantitative upper; it is not an experimentally demonstrated protocol.
